\begin{figure}[t]
  \caption{Model-implied departure effects by organizational resilience}
  \label{fig:model_departure_resilience}
  \centering
  \begin{minipage}[b]{0.325\textwidth}
    \centering
    \subcaption{Pull requests opened}
    \label{fig:model_resilience_opened}
    \includegraphics[width=\textwidth]{\ModelEventStudyDir/resilience/pc_score/norm/pull_request_opened.png}
  \end{minipage}
  \hfill
  \begin{minipage}[b]{0.325\textwidth}
    \centering
    \subcaption{Pull requests reviewed}
    \label{fig:model_resilience_reviewed}
    \includegraphics[width=\textwidth]{\ModelEventStudyDir/resilience/pc_score/norm/pull_request_reviewed.png}
  \end{minipage}
  \hfill
  \begin{minipage}[b]{0.325\textwidth}
    \centering
    \subcaption{Pull requests merged}
    \label{fig:model_resilience_merged}
    \includegraphics[width=\textwidth]{\ModelEventStudyDir/resilience/pc_score/norm/pull_request_merged.png}
  \end{minipage}

  \begin{minipage}{\textwidth}
    \textbf{Figure notes:} Each panel repeats Figure~\ref{fig:model_departure_effects} separately for organizations with high and low resilience. The shaded bands are the 95\% intervals of model-implied departure effects across \NModelDraws\ simulated datasets; solid lines are the actual event-study estimates. The merge-path outcomes appear in Appendix Figure~\ref{fig:model_departure_resilience_paths}.
  \end{minipage}
\todo[inline]{Improve figure notes.}
\end{figure}
